\section{THỬ NGHIỆM VÀ SO SÁNH KẾT QUẢ}
\subsection{Thiết lập mô hình kiểm chứng ResNet-18}
ResNet-18 là kiến trúc mạng tích chập có kết nối phần dư, giúp học phép biến đổi thông qua tổng giữa nhánh đặc trưng và nhánh đi tắt~\cite{resnet}. Trong bài tập lớn, lớp phân loại cuối được thay bằng ánh xạ từ 512 đặc trưng sang 43 lớp biển báo. Mô hình khởi tạo từ trọng số ImageNet và được tinh chỉnh trên tập ảnh gốc GTSRB.

Đầu vào ở khối xử lý ảnh có dạng $B\times3\times64\times64$ khi ghép batch. Trong hàm forward của ResNet-18, toàn bộ crop được nội suy song tuyến tính lên $224\times224$, có chống răng cưa, rồi chuẩn hóa theo trung bình $[0,485;0,456;0,406]$ và độ lệch chuẩn $[0,229;0,224;0,225]$. Chuẩn hóa được áp dụng đúng một lần. Đầu ra là ma trận $B\times43$.

\begin{table}[htbp]
\centering
\caption{Cấu hình huấn luyện ResNet-18 dùng để kiểm chứng.}
\begin{tabular}{P{6.1cm}P{8.0cm}}
\toprule
Thành phần & Thiết lập \\
\midrule
Kiến trúc & ResNet-18, lớp cuối 43 đầu ra \\
Khởi tạo & ImageNet, tinh chỉnh toàn mạng \\
Dữ liệu huấn luyện & Ảnh gốc; không huấn luyện riêng theo enhancer \\
Bộ tối ưu & AdamW, learning rate $10^{-4}$ \\
Weight decay / batch size & $10^{-4}$ / 32 \\
Số epoch tối đa & 20 \\
Dừng sớm & Patience 5, min delta 0,001 \\
Training seed & 11, 22, 33 \\
Chọn checkpoint & Macro-F1 validation clean tốt nhất \\
\bottomrule
\end{tabular}
\end{table}

Trong huấn luyện, các phép tăng cường nhẹ gồm xoay tối đa $10^\circ$, tịnh tiến tối đa 5\%, thay đổi nhẹ sáng/tương phản/bão hòa và xóa ngẫu nhiên một vùng nhỏ. Không lật ngang biển báo vì có thể thay đổi ý nghĩa hướng di chuyển. Cụm từ ``huấn luyện clean'' ở đây nghĩa là không huấn luyện bằng mô hình thiếu sáng Gamma--Poisson--Gaussian của thí nghiệm, không có nghĩa là loại bỏ mọi augmentation.

Ba checkpoint tốt nhất trong notebook nằm ở epoch 12, 10 và 3 tương ứng seed 11, 22 và 33; huấn luyện dừng ở epoch 17, 14 và 8. Sau bước này, trọng số được giữ cố định cho mọi phương pháp xử lý. Do vậy, chênh lệch dự đoán trong mỗi checkpoint phản ánh ảnh hưởng của đầu vào đã được biến đổi.

\subsection{Chỉ số đánh giá nhận dạng}
Accuracy là tỷ lệ mẫu được phân loại đúng:
\begin{equation}
\operatorname{Accuracy}=\frac{N_{\mathrm{đúng}}}{N_{\mathrm{tổng}}}.
\end{equation}
Với mỗi lớp $c$, Precision, Recall và F1 được tính từ số dự đoán đúng dương (TP), sai dương (FP) và sai âm (FN):
\begin{align}
P_c&=\frac{TP_c}{TP_c+FP_c},\qquad
R_c=\frac{TP_c}{TP_c+FN_c},\\
F1_c&=\frac{2TP_c}{2TP_c+FP_c+FN_c},\qquad
\operatorname{Macro\text{-}F1}=\frac{1}{43}\sum_{c=1}^{43}F1_c.
\end{align}
Mỗi lớp đóng góp như nhau vào Macro-F1, phù hợp với yêu cầu nhận dạng đầy đủ 43 loại biển báo. Weighted-F1 dùng số mẫu mỗi lớp làm trọng số và được lưu như metric bổ sung. Với tập con có đúng 10 ảnh/lớp, Weighted-F1 bằng Macro-F1 về mặt toán học; trên full-test có tần suất lớp khác nhau, hai chỉ số không nhất thiết bằng nhau.

Các bảng F1 trong báo cáo dùng Macro-F1 và ghi rõ đơn vị phần trăm. Chênh lệch giữa hai giá trị phần trăm được gọi là \textit{điểm phần trăm}, tránh nhầm với tỷ lệ cải thiện tương đối.

\subsection{Giao thức chọn tham số và tổng hợp số liệu}
Mỗi phương pháp có tham số được chọn trên 1.290 ảnh validation. Điểm lựa chọn là Macro-F1 trung bình qua ba mức suy giảm, ba noise seed và ba training seed. Gamma thử ba số mũ; CLAHE thử sáu cấu hình; MSR thử hai bộ tỉ lệ. HE không có tham số cần chọn trong cài đặt này. Ba bộ tăng cường sâu giữ cấu hình và trọng số công bố.

Sau khi chọn cấu hình tốt nhất của từng phương pháp, quy tắc chọn pipeline xét các ứng viên có Macro-F1 cách điểm tốt nhất dưới 0,005, tương đương 0,5 điểm phần trăm, rồi ưu tiên thời gian xử lý thấp hơn. Quy tắc này được xác định trước test. Trong notebook, diffusion có Macro-F1 validation khoảng 75,96\%, cao hơn CLAHE 72,02\%, do đó diffusion là phương pháp được chọn.

Số lượt đánh giá test của ma trận chính là:
\begin{equation}
3+3\times3\times3\times8=219.
\end{equation}
Ba lượt đầu là clean cho ba checkpoint. Phần còn lại là tích của training seed, mức suy giảm, noise seed và tám phương pháp tính cả Identity. Với tập con 430 ảnh, có 94.170 dự đoán. Notebook lưu kết quả kiểm chứng gồm hash checkpoint, độ phủ mã ảnh/nhãn và việc tính lại metric nhận dạng.

Khi tổng hợp, trước hết lấy trung bình qua noise seed trong từng training seed, sau đó tính trung bình và độ lệch chuẩn mẫu qua ba training seed:
\begin{equation}
\overline{m}_t=\frac{1}{3}\sum_{n=1}^{3}m_{t,n},\qquad
\mu=\frac{1}{3}\sum_{t=1}^{3}\overline{m}_t,\qquad
s=\sqrt{\frac{\sum_{t=1}^{3}(\overline{m}_t-\mu)^2}{2}}.
\end{equation}
Ba mức noise seed không được xem là ba lần huấn luyện độc lập. Các chỉ số ảnh không phụ thuộc training seed khi bộ tăng cường và ảnh đầu vào giữ nguyên. Độ lệch chuẩn qua training seed vì thế không phải thước đo biến thiên chất lượng ảnh giữa các mẫu.

\subsection{Nguồn kết quả và khả năng đối chiếu}
Phân tích chi tiết sử dụng notebook đã chạy \path{GTSRB_LLIE_Complete_Pipeline_Executed.ipynb}. Tệp này chứa hình ảnh, 219 bản ghi kết quả ResNet-18 trên 430 ảnh, bảng ablation và đầu ra kiểm chứng. Các bảng trong báo cáo được tổng hợp lại từ các bản ghi đó; hình minh họa được trích nguyên từ đầu ra notebook. Các cell được dùng gồm 7--9 (minh họa xử lý ảnh), 13--14 (validation và tham số đã chọn), 16--17 (kết quả test) và 20 (ablation), đánh số từ 0. Bảng full-test lấy từ README và được ghi nguồn riêng trong nội dung; bảng tổng hợp, tương quan và thông tin kiểm tra nguồn được lưu trong thư mục \texttt{report/evidence/}.

Bảng full-test 12.630 ảnh được lấy riêng từ README của dự án. Đây là thí nghiệm ResNet-18 khác với lần chạy notebook. Trong bản sao dự án dùng để soạn báo cáo chưa có thư mục kết quả full-test tương ứng, nên báo cáo không tuyên bố đã kiểm tra lại dự đoán từng ảnh của thí nghiệm này. Không ghép PSNR/SSIM của tập con với Macro-F1 full-test để tính tương quan.

\subsection{Chất lượng ảnh trên tập con 430 ảnh}
\begin{table}[htbp]
\centering
\caption{PSNR (dB) / SSIM của notebook ResNet-18 trên 430 ảnh. Mỗi ô ghi PSNR trước, SSIM sau.}
{\fontsize{11.5bp}{13.5bp}\selectfont
\begin{tabular}{lccc}
\toprule
Phương pháp & Tối nhẹ & Tối vừa & Tối nặng \\
\midrule
\tableinput{tables/quality}
\bottomrule
\end{tabular}}
\label{tab:quality}
\end{table}

Ở mức nhẹ, Gamma đạt PSNR 23,94 dB và SSIM 0,7011, cao hơn Identity ở cả hai chỉ số. Điều này phù hợp với tính chất của phép biến đổi Gamma khi suy giảm chưa quá mạnh: một mức nâng sáng vừa phải giúp giảm sai lệch cường độ mà chưa phải trả giá lớn về nhiễu.

Ở mức vừa, Zero-DCE có PSNR cao nhất trong bảng (20,08 dB), còn CLAHE có SSIM cao nhất (0,5568). Diffusion đạt PSNR 14,90 dB, thấp hơn Identity 15,37 dB, nhưng SSIM tăng từ 0,3410 lên 0,5410. Hai thước đo đã đưa ra những đánh giá khác nhau: sai lệch cường độ không đồng nhất với mức bảo toàn cấu trúc.

Ở mức nặng, Retinexformer có PSNR cao nhất (17,39 dB), trong khi diffusion có SSIM cao nhất (0,4720). MSR có PSNR thấp ở cả ba mức và SSIM giảm rõ khi nhiễu mạnh. Các kết quả này cho thấy không có phương pháp đứng đầu đồng thời ở mọi chỉ số và mọi điều kiện.

\subsection{Accuracy và Macro-F1 trên cùng tập con}
\begin{table}[htbp]
\centering
\caption{Accuracy trung bình $\pm$ độ lệch chuẩn mẫu (\%) của ResNet-18 trên 430 ảnh.}
{\fontsize{11.5bp}{13.5bp}\selectfont
\begin{tabular}{lccc}
\toprule
Phương pháp & Tối nhẹ & Tối vừa & Tối nặng \\
\midrule
\tableinput{tables/subset_acc}
\bottomrule
\end{tabular}}
\label{tab:acc}
\end{table}

\begin{table}[htbp]
\centering
\caption{Macro-F1 trung bình $\pm$ độ lệch chuẩn mẫu (\%) của ResNet-18 trên 430 ảnh.}
{\fontsize{11.5bp}{13.5bp}\selectfont
\begin{tabular}{lccc}
\toprule
Phương pháp & Tối nhẹ & Tối vừa & Tối nặng \\
\midrule
\tableinput{tables/subset_f1}
\bottomrule
\end{tabular}}
\label{tab:subsetf1}
\end{table}

Trên ảnh clean, Accuracy đạt $98,53\%\pm0,82$ và Macro-F1 đạt $98,50\%\pm0,83$. Khi giữ nguyên ảnh thiếu sáng, Macro-F1 giảm còn 90,38\%; 75,35\% và 54,81\% tương ứng mức nhẹ, vừa và nặng. Sự suy giảm này cho thấy mô hình đã học tốt trên ảnh gốc vẫn nhạy với thay đổi đầu vào.

Diffusion, được chọn bằng validation, đạt Macro-F1 92,01\%; 82,71\% và 63,72\%. So với Identity, mức tăng lần lượt xấp xỉ 1,63; 7,36 và 8,91 điểm phần trăm. Ở mức vừa, Accuracy tăng từ 73,05\% lên 82,95\%, cho thấy lợi ích không chỉ xuất hiện ở một cách tổng hợp F1.

CLAHE là phương án truyền thống nổi bật: Macro-F1 đạt 91,98\%; 80,14\% và 59,24\%. Riêng mức nhẹ, chênh lệch trung bình giữa CLAHE và diffusion chỉ khoảng 0,03 điểm phần trăm, nhỏ hơn nhiều độ lệch chuẩn giữa các seed. Không nên diễn giải khoảng cách rất nhỏ này thành bằng chứng ưu thế chắc chắn của một phương pháp.

\fig{subset_f1.png}{0.95}{Macro-F1 theo mức suy giảm trên tập con 430 ảnh của notebook ResNet-18. Các đường thể hiện xu hướng của cùng một bộ nhận dạng khi thay đổi phép xử lý đầu vào.}{fig:subsetf1}

\subsection{Đối chiếu trên toàn bộ 12.630 ảnh test}
\begin{table}[htbp]
\centering
\caption{Macro-F1 (\%) của ResNet-18 trên full-test, theo bảng đã công bố trong README. Đây là lần chạy riêng với notebook 430 ảnh.}
{\fontsize{11.5bp}{13.5bp}\selectfont
\begin{tabular}{lccc}
\toprule
Phương pháp & Tối nhẹ & Tối vừa & Tối nặng \\
\midrule
\tableinput{tables/fulltest}
\bottomrule
\end{tabular}}
\label{tab:full}
\end{table}

Theo README, full-test clean có Accuracy $99,15\%\pm0,30$ và Macro-F1 $98,90\%\pm0,40$. Pipeline diffusion đã chọn bằng validation đạt Macro-F1 93,45\%; 83,87\% và 64,88\% ở ba mức. Mức tăng so với Identity tương ứng là 1,45; 5,83 và 5,18 điểm phần trăm.

CLAHE có Macro-F1 mức nhẹ cao nhất trong bảng full-test (94,35\%). Tuy nhiên, đây là quan sát trên test, không phải căn cứ để thay phương pháp đã chọn trước đó. Ở mức vừa và nặng, diffusion có trung bình cao hơn CLAHE, nhưng kết luận về ý nghĩa thống kê cần phân tích chênh lệch ghép cặp từ dự đoán gốc.

Kết quả full-test củng cố nhận xét rằng thiếu sáng nặng vẫn là trường hợp khó và tăng sáng không phục hồi hoàn toàn hiệu năng clean. Các giá trị khác với bảng tập con là điều bình thường khi phạm vi mẫu và lần chạy thay đổi. Báo cáo sử dụng bảng này để đối chiếu xu hướng, không xem hai bảng là cùng một thí nghiệm.
